\subsection{Purpose}

Patch and vulnerability-management procedures identify, evaluate, prioritize, remediate, and monitor weaknesses in systems, applications, devices, networks, cloud resources, and other technology assets.

The purpose of these procedures is to reduce the likelihood that vulnerabilities will be exploited and to ensure that security updates are applied in a timely and controlled manner.

\subsection{Asset Coverage}

Patch and vulnerability management \must{} apply to all supported technology resources, including:
\begin{itemize}
    \item Servers, workstations, laptops, and mobile devices
    \item Operating systems, applications, databases, and middleware
    \item Network devices, firewalls, wireless systems, and VPN equipment
    \item Virtual machines, containers, hypervisors, and cloud resources
    \item Security appliances and management platforms
    \item Firmware, drivers, libraries, and third-party components
    \item Internet-facing systems and externally hosted services
\end{itemize}

Systems and applications \must{} be included in the asset inventory so that their ownership, software versions, support status, and patch requirements can be tracked.

Unsupported or end-of-life systems \must{} be identified and documented. Such systems \must{} be replaced, upgraded, isolated, or retired as soon as reasonably practicable.

\subsection{Vulnerability Identification}

The organization \must{} use appropriate sources to identify vulnerabilities and security updates affecting its systems.

Vulnerability information may be obtained from:
\begin{itemize}
    \item Software and hardware vendors
    \item Operating system and application security advisories
    \item Vulnerability databases
    \item Security monitoring services
    \item Internal vulnerability scans
    \item Penetration tests and security assessments
    \item Incident investigations
    \item Managed service providers and security partners
\end{itemize}

Administrators and system owners \must{} monitor for vulnerabilities affecting systems under their responsibility.

Vulnerability scans \should{} be performed periodically on internal and externally accessible systems. Internet-facing systems \should{} be scanned more frequently based on their exposure and risk.

\subsection{Vulnerability Assessment}

Identified vulnerabilities \must{} be evaluated according to their severity, exploitability, affected assets, exposure, data sensitivity, and potential business impact.

The assessment \should{} consider:
\begin{itemize}
    \item The severity assigned by the vendor or vulnerability source
    \item Whether the vulnerability is actively being exploited
    \item Whether public exploit code is available
    \item Whether the affected system is accessible from the Internet
    \item The sensitivity and criticality of the affected system
    \item The availability of a patch or other remediation
    \item Existing compensating controls
    \item The operational risk of applying the remediation
\end{itemize}

Vulnerabilities \must{} be prioritized based on risk rather than severity alone.

\subsection{Patch Priorities}

Unless a documented exception is approved, patches and vulnerability remediations \should{} be addressed according to the following priorities:
\begin{itemize}
    \item Critical vulnerabilities affecting Internet-facing or mission-critical systems \should{} be remediated as soon as possible.
    \item High-risk vulnerabilities \should{} be remediated within a defined period based on organizational risk tolerance.
    \item Medium-risk vulnerabilities \should{} be remediated during the next appropriate maintenance cycle.
    \item Low-risk vulnerabilities \should{} be addressed through routine maintenance or other planned remediation.
\end{itemize}

Vulnerabilities known to be actively exploited \must{} receive urgent attention regardless of their original severity rating.

The organization \must{} define target remediation timeframes for each risk level. Missed remediation timeframes \must{} be documented and escalated to the system owner or responsible manager.

\subsection{Patch Testing}

Patches \should{} be tested before deployment when testing is reasonably practicable and the risk of testing is lower than the risk of immediate deployment.

Testing \should{} verify:
\begin{itemize}
    \item The system remains operational
    \item Required applications and services continue to function
    \item Security controls remain enabled
    \item Business-critical integrations continue to operate
    \item Backups and recovery procedures remain available
\end{itemize}

Testing may be limited or skipped when a vulnerability is actively exploited, the affected system is highly exposed, or the vendor identifies the update as an urgent security fix.

The reason for bypassing normal testing \must{} be documented.

\subsection{Patch Deployment}

Patches \must{} be obtained from approved and trusted sources.

Patches \must{} be applied using approved procedures and tools.

Administrators \must{} verify, where technically practicable, that:
\begin{itemize}
    \item The patch is intended for the affected system
    \item The patch was obtained from an approved source
    \item The patch installation completed successfully
    \item The system restarted or completed any required post-installation tasks
    \item The vulnerability or outdated version is no longer present
\end{itemize}

Automatic updates \should{} be enabled where they are supported and appropriate for the system.

Security updates \must{} not be permanently deferred merely for convenience or to avoid routine maintenance.

\subsection{Emergency Remediation}

Emergency remediation may include expedited patching, temporary isolation, disabling an affected service, restricting network access, applying a vendor workaround, or taking the affected system offline.

Emergency remediation \must{} be considered when:
\begin{itemize}
    \item A vulnerability is actively exploited
    \item A critical system is exposed to an untrusted network
    \item Sensitive information may be accessed through the vulnerability
    \item Exploit code is publicly available
    \item No acceptable compensating control is available
\end{itemize}

Emergency actions \must{} be documented and reviewed after implementation.

\subsection{Verification and Reporting}

Patch and vulnerability remediation \must{} be verified after deployment.

Verification may include:
\begin{itemize}
    \item Reviewing patch-management reports
    \item Checking installed software and firmware versions
    \item Performing follow-up vulnerability scans
    \item Reviewing system logs
    \item Confirming that required services remain operational
    \item Obtaining confirmation from the system owner
\end{itemize}

Patch and vulnerability-management records \must{} identify the affected asset, vulnerability or update, risk rating, remediation action, responsible person, date identified, date remediated, and verification result.

Management \should{} receive periodic reports showing outstanding vulnerabilities, overdue patches, unsupported systems, and significant remediation risks.

\subsection{Exceptions and Compensating Controls}

A patch or vulnerability that cannot be remediated within the required timeframe \must{} have a documented exception.

The exception record \must{} include:
\begin{itemize}
    \item The affected system and vulnerability
    \item The reason remediation cannot be completed
    \item The risks involved
    \item Compensating controls
    \item The responsible system owner
    \item The approving authority
    \item The planned remediation date
    \item The expiration or review date
\end{itemize}

Compensating controls may include network isolation, access restrictions, application allowlisting, additional monitoring, disabling vulnerable functionality, or increased logging.

Exceptions \must{} be reviewed periodically and closed when remediation is completed or the affected system is retired.